\documentclass[12pt]{article}
\usepackage[utf8]{inputenc}
\usepackage[T1]{fontenc}
\usepackage{hyperref}
\title{Titan Language Resources for K-12 Curriculum Development}
\author{Compiled by Claude Sonnet 5 (claude-sonnet-5) for Andr\'as Kornai}
\date{}
\begin{document}
\maketitle

\section{Introduction}
Titan (ISO 639-3 ttv; Glottocode tita1241) is an Austronesian language of the East Manus group. It is spoken on the southeastern part of Manus Island and neighbouring islands in Papua New Guinea. Sources give about 4,000 to 5,100 speakers. Searches found a sketch grammar and archival collections. They did not find a published stand-alone dictionary, nor any school materials in Titan.

\section{Primary Educational Resources}
No primers, textbooks or school materials in Titan turned up in the searches. Any curriculum work would have to start from the linguistic documentation listed below.

\section{Linguistic and Reference Materials}
\begin{itemize}
\item Titan language overview (Wikipedia): \url{https://en.wikipedia.org/wiki/Titan_language}
\item Manus languages overview (Wikipedia): \url{https://en.wikipedia.org/wiki/Manus_languages}
\item Titan Materials, a research data record that includes dictionary text files: \url{https://researchdata.edu.au/titan-materials/1425841}
\end{itemize}
The search summary names Claire Bowern's \emph{Sivisa Titan: Sketch Grammar, Texts, Vocabulary} (Oceanic Linguistics Special Publication 38, University of Hawai'i Press, 2011). I did not see a direct URL for the book in the results, so none is given here.

\section{Religious and Community Resources}
Nothing relevant was found in the searches.

\section{Digital and Audio Resources}
\begin{itemize}
\item PARADISEC Titan grammar master file: \url{https://catalog.paradisec.org.au/repository/CB4/TitanGrammarMaster}
\item PARADISEC collection CB4 (Titan materials): \url{https://catalog.nabu-prod.paradisec.org.au/collections/CB4}
\end{itemize}
The search summary also mentions Robert Blust's Titan audio and written field materials at Kaipuleohone. No URL for them appeared in the results.

\section{Curriculum Development Context}
The language is stable and is used as a first language across the community. The searches found no evidence that it is taught in schools. The existing documentation (sketch grammar, vocabulary, texts, archival recordings) is the likely basis for developing an orthography-based primer and other materials. Little else was found.
\end{document}